%% USPSC-Cap5-Resultados.tex
% ----------------------------------------------------------
% Resultados e Discussão
% ----------------------------------------------------------
%
% ATENÇÃO: capítulo estruturado antes da coleta.
% Todos os quadros e tabelas estão montados e rotulados, com linhas de
% exemplo marcadas como [preencher]. Ao obter os dados, basta substituir o
% conteúdo das células e redigir os parágrafos de análise indicados.
%
% Convenção adotada:
%   Rodada 1 = avaliação inicial, ANTES das melhorias  (participantes P1..P3)
%   Rodada 2 = reavaliação, APÓS as melhorias          (participantes P4..P6)
% ----------------------------------------------------------
\chapter{Resultados e Discussão}
\label{cap:resultados}

\section{Considerações Iniciais}
\label{sec:res_consideracoes}

Este capítulo apresenta os resultados obtidos nas quatro fases descritas no
Capítulo~\ref{cap:metodologia}. A Seção~\ref{sec:res_heuristica} relata os achados
da avaliação heurística e as correções dela decorrentes. A
Seção~\ref{sec:res_rodada1} apresenta os resultados da avaliação inicial com
usuários. A Seção~\ref{sec:res_correcoes} descreve as melhorias implementadas em
resposta a esses achados. A Seção~\ref{sec:res_rodada2} apresenta os resultados da
reavaliação. A Seção~\ref{sec:res_matriz} consolida a matriz de rastreabilidade e
a Seção~\ref{sec:res_discussao} discute os achados à luz da questão de pesquisa.

\section{Resultados da Avaliação Heurística}
\label{sec:res_heuristica}

% >>> [PREENCHER] Parágrafo introdutório: número de avaliadores, perfil de cada
% um, telas inspecionadas, duração da inspeção e forma de consolidação dos
% relatos individuais.

\begin{quadro}[htb]
\centering
\caption{Problemas identificados na avaliação heurística}
\label{quad:res_heuristica}
\footnotesize
\begin{tabular}{p{1.0cm} p{3.0cm} p{4.4cm} p{1.2cm} p{1.2cm} p{1.4cm}}
\toprule
\textbf{ID} & \textbf{Tela} & \textbf{Problema} & \textbf{Heur.} &
\textbf{Sev.} & \textbf{Corrigido} \\
\midrule
AH01 & \textit{[preencher]} & \textit{[preencher]} & --- & --- & --- \\
AH02 & \textit{[preencher]} & \textit{[preencher]} & --- & --- & --- \\
AH03 & \textit{[preencher]} & \textit{[preencher]} & --- & --- & --- \\
\bottomrule
\end{tabular}
\\[4pt]
\begin{flushleft}
Fonte: Elaborado pelo autor.
\end{flushleft}
\end{quadro}

% >>> [PREENCHER] Análise: distribuição dos problemas por heurística e por
% severidade; quais foram corrigidos antes da Fase 1 e quais foram adiados,
% com a justificativa do adiamento.

\section{Rodada 1 --- Avaliação Inicial}
\label{sec:res_rodada1}

\subsection{Caracterização dos participantes e das condições}

\begin{quadro}[htb]
\centering
\caption{Perfil dos participantes da Rodada 1}
\label{quad:res_perfil_r1}
\footnotesize
\begin{tabular}{p{1.0cm} p{1.0cm} p{2.8cm} p{2.6cm} p{2.4cm} p{2.0cm}}
\toprule
\textbf{ID} & \textbf{Idade} & \textbf{Ocupação} & \textbf{Proto-persona
correspondente} & \textbf{Uso de apps de percurso} & \textbf{Esforço percebido
(0--10)} \\
\midrule
P1 & --- & \textit{[preencher]} & --- & --- & --- \\
P2 & --- & \textit{[preencher]} & --- & --- & --- \\
P3 & --- & \textit{[preencher]} & --- & --- & --- \\
\bottomrule
\end{tabular}
\\[4pt]
\begin{flushleft}
Fonte: Elaborado pelo autor.
\end{flushleft}
\end{quadro}

% >>> [PREENCHER] Registrar: data e horário das sessões, condições
% meteorológicas, trilha utilizada, número de checkpoints e versão do
% aplicativo. Comentar a aderência entre proto-personas planejadas e
% participantes efetivamente recrutados, justificando eventuais desvios.

\subsection{Desempenho nas tarefas}

\begin{quadro}[htb]
\centering
\caption{Conclusão das tarefas na Rodada 1}
\label{quad:res_tarefas_r1}
\footnotesize
\begin{tabular}{c p{3.2cm} c c c p{2.6cm}}
\toprule
\textbf{Tarefa} & \textbf{Descrição resumida} & \textbf{P1} & \textbf{P2} &
\textbf{P3} & \textbf{Tempo médio} \\
\midrule
T1 & Descoberta e escolha da trilha & --- & --- & --- & --- \\
T2 & Leitura da tela de preparação & --- & --- & --- & --- \\
T3 & Deslocamento ao \textit{checkpoint} & --- & --- & --- & --- \\
T4 & Registro por QR Code & --- & --- & --- & --- \\
T5 & Leitura do estado da atividade & --- & --- & --- & --- \\
T6 & Suspensão e retomada & --- & --- & --- & --- \\
T7 & Conclusão da trilha & --- & --- & --- & --- \\
\bottomrule
\end{tabular}
\\[4pt]
\begin{flushleft}
Legenda: A = autônoma; D$n$ = com auxílio, grau $n$; F = falha.\\
Fonte: Elaborado pelo autor.
\end{flushleft}
\end{quadro}

\subsection{Episódios de atrito}

% >>> [PREENCHER] Para cada episódio, redigir um parágrafo curto contendo:
% a situação, a verbalização literal do participante (entre aspas), a
% heurística tocada e a origem atribuída. Agrupar por tarefa.

\begin{quadro}[htb]
\centering
\caption{Episódios de atrito registrados na Rodada 1}
\label{quad:res_atritos_r1}
\footnotesize
\begin{tabular}{p{1.0cm} p{1.0cm} p{4.0cm} p{1.2cm} p{1.8cm} p{1.2cm}}
\toprule
\textbf{ID} & \textbf{Part.} & \textbf{Descrição} & \textbf{Heur.} &
\textbf{Origem} & \textbf{Sev.} \\
\midrule
A01 & --- & \textit{[preencher]} & --- & --- & --- \\
A02 & --- & \textit{[preencher]} & --- & --- & --- \\
A03 & --- & \textit{[preencher]} & --- & --- & --- \\
\bottomrule
\end{tabular}
\\[4pt]
\begin{flushleft}
Origem: I = interface; D = desafio proposto; F = infraestrutura;
T = dívida técnica.\\
Fonte: Elaborado pelo autor.
\end{flushleft}
\end{quadro}

\subsection{Distribuição dos atritos por origem}

% >>> [PREENCHER] Esta é a subseção que responde diretamente à questão de
% pesquisa. Apresentar a contagem de episódios por categoria de origem e
% discutir os casos de atribuição ambígua, isto é, aqueles em que a
% verbalização do participante e a interpretação do observador divergiram.

\begin{quadro}[htb]
\centering
\caption{Distribuição dos episódios de atrito por origem --- Rodada 1}
\label{quad:res_origem_r1}
\begin{tabular}{p{5.4cm} c c}
\toprule
\textbf{Origem do atrito} & \textbf{Episódios} & \textbf{Proporção} \\
\midrule
Interface & --- & --- \\
Desafio proposto & --- & --- \\
Infraestrutura & --- & --- \\
Dívida técnica & --- & --- \\
\midrule
\textbf{Total} & --- & 100\,\% \\
\bottomrule
\end{tabular}
\\[4pt]
\begin{flushleft}
Fonte: Elaborado pelo autor.
\end{flushleft}
\end{quadro}

\subsection{Achados da entrevista semiestruturada}

% >>> [PREENCHER] Organizar por bloco do roteiro:
%   (a) percepção geral e compreensão da proposta;
%   (b) atribuição da dificuldade pelo próprio participante;
%   (c) compreensão dos indicadores de estado;
%   (d) reconhecimento de ícones --- reportar quantos e quais ícones foram
%       corretamente identificados por participante;
%   (e) modelo mental da validação por QR Code e percepção de confiança;
%   (f) priorização espontânea de melhorias.

\begin{quadro}[htb]
\centering
\caption{Reconhecimento de ícones --- Rodada 1}
\label{quad:res_icones_r1}
\begin{tabular}{p{5.0cm} c c c}
\toprule
\textbf{Elemento} & \textbf{P1} & \textbf{P2} & \textbf{P3} \\
\midrule
Recentralizar mapa & --- & --- & --- \\
Bússola & --- & --- & --- \\
Marcador de posição do usuário & --- & --- & --- \\
Marcadores coloridos de \textit{checkpoint} & --- & --- & --- \\
Botões circulares de ação & --- & --- & --- \\
\bottomrule
\end{tabular}
\\[4pt]
\begin{flushleft}
Legenda: C = reconhecido corretamente; P = parcialmente; N = não reconhecido.\\
Fonte: Elaborado pelo autor.
\end{flushleft}
\end{quadro}

\subsection{Resultados do instrumento TAM}

Conforme justificado na Seção~\ref{sec:instrumento_tam}, os escores da Rodada 1 são
reportados de forma descritiva e individual, sem agregação inferencial.

\begin{quadro}[htb]
\centering
\caption{Escores do instrumento TAM --- Rodada 1}
\label{quad:res_tam_r1}
\begin{tabular}{p{2.0cm} c c c c}
\toprule
\textbf{Item} & \textbf{P1} & \textbf{P2} & \textbf{P3} & \textbf{Mediana} \\
\midrule
PEOU1 & --- & --- & --- & --- \\
PEOU2 & --- & --- & --- & --- \\
PEOU3 & --- & --- & --- & --- \\
PEOU4 & --- & --- & --- & --- \\
\midrule
PU1 & --- & --- & --- & --- \\
PU2 & --- & --- & --- & --- \\
PU3 & --- & --- & --- & --- \\
PU4 & --- & --- & --- & --- \\
\midrule
ITU1 & --- & --- & --- & --- \\
ITU2 & --- & --- & --- & --- \\
\bottomrule
\end{tabular}
\\[4pt]
\begin{flushleft}
Escala Likert de sete pontos.\\
Fonte: Elaborado pelo autor.
\end{flushleft}
\end{quadro}

\subsection{Validação semântica do instrumento}

% >>> [PREENCHER] Relatar, item a item, se houve dúvida de interpretação e
% qual reformulação o participante produziu ao explicar o enunciado. Indicar
% quais itens demandam revisão de redação antes da Rodada 2 e registrar a
% redação revista.

\section{Correções Implementadas}
\label{sec:res_correcoes}

% >>> [PREENCHER] Descrever cada correção, vinculando-a ao episódio de atrito
% ou ao problema heurístico que a motivou. Incluir capturas de tela do antes e
% do depois quando a alteração for visual.

\begin{quadro}[htb]
\centering
\caption{Correções implementadas entre as rodadas}
\label{quad:res_correcoes}
\footnotesize
\begin{tabular}{p{1.0cm} p{4.6cm} p{4.0cm} p{2.4cm}}
\toprule
\textbf{ID} & \textbf{Correção} & \textbf{Achado que a motivou} &
\textbf{Camada alterada} \\
\midrule
C01 & \textit{[preencher]} & --- & --- \\
C02 & \textit{[preencher]} & --- & --- \\
C03 & \textit{[preencher]} & --- & --- \\
\bottomrule
\end{tabular}
\\[4pt]
\begin{flushleft}
Fonte: Elaborado pelo autor.
\end{flushleft}
\end{quadro}

% >>> [PREENCHER] Registrar também os achados que NÃO geraram correção, com a
% justificativa: atritos classificados como decorrentes do desafio proposto,
% correções de custo incompatível com o prazo, ou problemas cuja solução
% depende de infraestrutura externa ao projeto.

\section{Rodada 2 --- Reavaliação após as Melhorias}
\label{sec:res_rodada2}

% >>> [PREENCHER] Replicar aqui a mesma estrutura de quadros da Rodada 1,
% com os participantes P4 a P6. Antes disso, registrar explicitamente que as
% condições foram replicadas: mesma trilha, mesma faixa de horário, mesmo
% dispositivo, e informar a nova versão do aplicativo.

\subsection{Caracterização dos participantes e das condições}

\begin{quadro}[htb]
\centering
\caption{Perfil dos participantes da Rodada 2}
\label{quad:res_perfil_r2}
\footnotesize
\begin{tabular}{p{1.0cm} p{1.0cm} p{2.8cm} p{2.6cm} p{2.4cm} p{2.0cm}}
\toprule
\textbf{ID} & \textbf{Idade} & \textbf{Ocupação} & \textbf{Proto-persona
correspondente} & \textbf{Uso de apps de percurso} & \textbf{Esforço percebido
(0--10)} \\
\midrule
P4 & --- & \textit{[preencher]} & --- & --- & --- \\
P5 & --- & \textit{[preencher]} & --- & --- & --- \\
P6 & --- & \textit{[preencher]} & --- & --- & --- \\
\bottomrule
\end{tabular}
\\[4pt]
\begin{flushleft}
Fonte: Elaborado pelo autor.
\end{flushleft}
\end{quadro}

\subsection{Desempenho nas tarefas}

\begin{quadro}[htb]
\centering
\caption{Conclusão das tarefas na Rodada 2}
\label{quad:res_tarefas_r2}
\footnotesize
\begin{tabular}{c p{3.2cm} c c c p{2.6cm}}
\toprule
\textbf{Tarefa} & \textbf{Descrição resumida} & \textbf{P4} & \textbf{P5} &
\textbf{P6} & \textbf{Tempo médio} \\
\midrule
T1 & Descoberta e escolha da trilha & --- & --- & --- & --- \\
T2 & Leitura da tela de preparação & --- & --- & --- & --- \\
T3 & Deslocamento ao \textit{checkpoint} & --- & --- & --- & --- \\
T4 & Registro por QR Code & --- & --- & --- & --- \\
T5 & Leitura do estado da atividade & --- & --- & --- & --- \\
T6 & Suspensão e retomada & --- & --- & --- & --- \\
T7 & Conclusão da trilha & --- & --- & --- & --- \\
\bottomrule
\end{tabular}
\\[4pt]
\begin{flushleft}
Fonte: Elaborado pelo autor.
\end{flushleft}
\end{quadro}

\subsection{Episódios de atrito e escores do instrumento}

% >>> [PREENCHER] Reproduzir os quadros de atritos, de distribuição por origem,
% de reconhecimento de ícones e de escores TAM, com os rótulos
% quad:res_atritos_r2, quad:res_origem_r2, quad:res_icones_r2 e
% quad:res_tam_r2.

\section{Comparação entre as Rodadas}
\label{sec:res_comparacao}

% >>> [PREENCHER] ATENÇÃO À NATUREZA DA AFIRMAÇÃO. Com três participantes por
% rodada, não cabe teste de hipótese nem afirmação de significância. A
% comparação legítima é qualitativa e rastreável, na forma:
% "o atrito A0x, observado em 3 dos 3 participantes da Rodada 1, não se
% manifestou em nenhum dos 3 participantes da Rodada 2, após a correção C0y".
% Evitar expressões como "houve melhora estatisticamente relevante".

\begin{quadro}[htb]
\centering
\caption{Persistência dos atritos entre as rodadas}
\label{quad:res_comparacao}
\footnotesize
\begin{tabular}{p{1.0cm} p{3.8cm} c c p{2.4cm} p{2.0cm}}
\toprule
\textbf{ID} & \textbf{Atrito} & \textbf{R1} & \textbf{R2} &
\textbf{Correção aplicada} & \textbf{Situação} \\
\midrule
A01 & \textit{[preencher]} & --- & --- & --- & --- \\
A02 & \textit{[preencher]} & --- & --- & --- & --- \\
\bottomrule
\end{tabular}
\\[4pt]
\begin{flushleft}
Situação: R = resolvido; P = persistente; N = novo; NA = não aplicável.\\
Fonte: Elaborado pelo autor.
\end{flushleft}
\end{quadro}

% >>> [PREENCHER] Comentar também os atritos NOVOS surgidos na Rodada 2. Sua
% ocorrência não invalida o trabalho: é achado esperado em projeto iterativo e
% deve ser reportado com a mesma transparência.

\section{Matriz de Rastreabilidade Consolidada}
\label{sec:res_matriz}

A matriz apresentada no Quadro~\ref{quad:res_matriz} materializa a contribuição
central enunciada na Seção~\ref{sec:problema}, ao vincular cada problema observado
à heurística violada, à decisão de engenharia que o originou, ao construto do
instrumento de aceitação potencialmente impactado e à correção proposta ou
aplicada.

\begin{quadro}[htb]
\centering
\caption{Matriz de rastreabilidade entre atrito, heurística, decisão técnica e construto}
\label{quad:res_matriz}
\scriptsize
\begin{tabular}{p{0.8cm} p{2.4cm} p{1.0cm} p{2.4cm} p{1.2cm} p{1.2cm} p{2.2cm}}
\toprule
\textbf{ID} & \textbf{Problema} & \textbf{Heur.} & \textbf{Decisão de origem} &
\textbf{Origem} & \textbf{Construto} & \textbf{Correção} \\
\midrule
P01 & --- & --- & --- & --- & --- & --- \\
P02 & --- & --- & --- & --- & --- & --- \\
P03 & --- & --- & --- & --- & --- & --- \\
\bottomrule
\end{tabular}
\\[4pt]
\begin{flushleft}
Fonte: Elaborado pelo autor.
\end{flushleft}
\end{quadro}

\section{Discussão}
\label{sec:res_discussao}

% >>> [PREENCHER] Organizar a discussão em torno de quatro eixos:
%
% (a) RESPOSTA À QUESTÃO DE PESQUISA. Retomar a pergunta da Seção 1.2 e
%     responder com base na distribuição dos atritos por origem: em que medida
%     a interface atuou como facilitadora, e a classificação proposta mostrou-se
%     operacionalizável na prática?
%
% (b) VERIFICAÇÃO DAS HIPÓTESES DO QUADRO 3.11. As limitações declaradas no
%     Capítulo 3 --- servimento de imagens, dependência de conectividade, tela
%     de perfil não integrada --- manifestaram-se como atritos observáveis?
%     Confirmar ou refutar cada uma, item a item.
%
% (c) CONVERGÊNCIA ENTRE OS INSTRUMENTOS. Os atritos verbalizados encontraram
%     eco nos escores de PEOU? Onde houve divergência entre o observado e o
%     declarado, discutir a possível causa.
%
% (d) CONFRONTO COM A LITERATURA. Comparar os achados com os trabalhos
%     relacionados levantados na Seção 2.8.

\section{Considerações Finais}
\label{sec:res_finais}

% >>> [PREENCHER] Sintetizar os principais achados em um parágrafo e anunciar
% o capítulo de conclusão.
